% preamble-exam-zh.tex — 共享 preamble
% 用于所有 exam-zh 模板
%
% 样式策略：
%   屏幕 → 语义彩色（蓝=关键想法，红=易错，绿=路线，灰=提示/教师备注）
%   打印 → 自动降级为黑白（通过 print mode 或 monochrome 选项）

\documentclass{exam-zh}

% ---------- 基础配置 ----------
\examsetup{
  page/size = a4paper,
  page/show-head = false,
  page/show-foot = false,
  sealline/show = false,
  font = times,
  math-font = xits,
}

\usepackage{tcolorbox}
\usepackage{needspace}
\tcbuselibrary{skins,breakable}

% ---------- 打印/屏幕颜色切换 ----------
% 用法：\documentclass[monochrome]{exam-zh} 即可切换为纯黑白
% 注意：不要使用 \if@xxx 放在 makeatother 外部；这里改成普通条件名。
\newif\ifeduprintcolor
\eduprintcolortrue

\makeatletter
\@ifclasswith{exam-zh}{monochrome}{\eduprintcolorfalse}{}
\makeatother

% ---------- 语义颜色定义 ----------
% 屏幕：有色彩区分；打印：降级为灰度
\ifeduprintcolor
  % --- 屏幕彩色模式 ---
  \definecolor{edu-blue}{RGB}{47, 82, 122}       % 关键想法
  \definecolor{edu-blue-bg}{RGB}{243, 246, 252}
  \definecolor{edu-red}{RGB}{180, 40, 40}         % 易错提醒
  \definecolor{edu-red-bg}{RGB}{255, 245, 240}
  \definecolor{edu-green}{RGB}{34, 100, 60}       % 路线图
  \definecolor{edu-green-bg}{RGB}{242, 250, 245}
  \definecolor{edu-orange}{RGB}{160, 90, 0}       % 训练目标
  \definecolor{edu-orange-bg}{RGB}{255, 248, 235}
  \definecolor{edu-gray}{RGB}{100, 100, 100}      % 提示/教师备注
  \definecolor{edu-gray-bg}{RGB}{248, 248, 248}
  \definecolor{edu-purple}{RGB}{90, 50, 120}      % 思考题
  \definecolor{edu-purple-bg}{RGB}{248, 244, 252}
\else
  % --- 打印黑白模式 ---
  \definecolor{edu-blue}{gray}{0.15}
  \definecolor{edu-blue-bg}{gray}{0.96}
  \definecolor{edu-red}{gray}{0.25}
  \definecolor{edu-red-bg}{gray}{0.97}
  \definecolor{edu-green}{gray}{0.20}
  \definecolor{edu-green-bg}{gray}{0.97}
  \definecolor{edu-orange}{gray}{0.25}
  \definecolor{edu-orange-bg}{gray}{0.98}
  \definecolor{edu-gray}{gray}{0.40}
  \definecolor{edu-gray-bg}{gray}{0.97}
  \definecolor{edu-purple}{gray}{0.30}
  \definecolor{edu-purple-bg}{gray}{0.98}
\fi

% ---------- 自定义教学环境 ----------
% 共性：enhanced + breakable（跨页不断裂）

% 原题卡片 — 强边框，突出显示
\newtcolorbox{problemcard}{
  enhanced, breakable,
  colback=edu-blue-bg,
  colframe=edu-blue,
  boxrule=1.2pt,
  arc=0pt,
  left=3mm, right=3mm, top=3mm, bottom=3mm,
}

% 解题路线图 — 左边线 + 浅底
\newtcolorbox{routebox}{
  enhanced, breakable,
  colback=edu-green-bg,
  colframe=edu-green!30,
  boxrule=0pt,
  borderline west={2.5pt}{0pt}{edu-green},
  left=3mm, right=3mm, top=3mm, bottom=3mm,
}

% 关键想法 — 蓝色左边线
\newtcolorbox{keyideabox}{
  enhanced, breakable,
  colback=edu-blue-bg,
  colframe=edu-blue!30,
  boxrule=0pt,
  borderline west={2.5pt}{0pt}{edu-blue},
  left=3mm, right=3mm, top=3mm, bottom=3mm,
}

% 易错提醒 — 红色左边线
\newtcolorbox{mistakebox}{
  enhanced, breakable,
  colback=edu-red-bg,
  colframe=edu-red!20,
  boxrule=0pt,
  borderline west={3pt}{0pt}{edu-red},
  left=3mm, right=3mm, top=2.5mm, bottom=2.5mm,
}

% 提示 — 灰色虚线左边线
\newtcolorbox{hintbox}{
  enhanced, breakable,
  colback=edu-gray-bg,
  colframe=edu-gray!20,
  boxrule=0pt,
  borderline west={2pt}{0pt}{edu-gray, dashed},
  left=3mm, right=3mm, top=2mm, bottom=2mm,
}

% 思考题 — 紫色虚线左边线
\newtcolorbox{thinkbox}{
  enhanced, breakable,
  colback=edu-purple-bg,
  colframe=edu-purple!20,
  boxrule=0pt,
  borderline west={2pt}{0pt}{edu-purple, dashed},
  left=2.5mm, right=2.5mm, top=2.5mm, bottom=2.5mm,
}

% 教师备注 — 灰色左边线，小字
\newtcolorbox{teachernote}{
  enhanced, breakable,
  colback=edu-gray-bg,
  colframe=edu-gray!15,
  boxrule=0pt,
  borderline west={2pt}{0pt}{edu-gray!60},
  left=2.5mm, right=2.5mm, top=2.5mm, bottom=2.5mm,
  fontupper=\small,
  colupper=black!60,
}

% 训练目标 — 橙色左边线，小字
\newtcolorbox{traininggoal}{
  enhanced, breakable,
  colback=edu-orange-bg,
  colframe=edu-orange!15,
  boxrule=0pt,
  borderline west={1.5pt}{0pt}{edu-orange},
  left=2mm, right=2mm, top=1mm, bottom=1mm,
  fontupper=\small,
}

% 答题区 — 无边框，纯留白
\newcommand{\answerarea}[1][40mm]{%
  \vspace{2mm}%
  \begin{tcolorbox}[
    enhanced, breakable,
    colback=white,
    colframe=white,
    boxrule=0pt,
    height=#1,
    valign=top,
    bottom=0mm,
    after skip=0mm,
  ]
  \end{tcolorbox}%
}

% ---------- 字体设置 ----------
% exam-zh (ctexbook) already sets CJK fonts via ctex.
% Only override if specific fonts are needed.
% macOS: Songti SC, STHeiti, STFangsong
% Linux: SimSun, SimHei, FangSong

% ---------- 页面设置 ----------
% exam-zh (ctexbook) already loads geometry.
% Override margins if needed:
% \geometry{top=16mm, bottom=16mm, left=14mm, right=14mm}

\examsetup{
  question/show-answer = true,
  fillin/show-answer = true,
  solution/show-solution = show-stay,
}

% ---------- 教师版解答样式 ----------
\newtcolorbox{solutionblock}{
  enhanced, breakable,
  colback=edu-blue-bg,
  colframe=edu-blue!25,
  boxrule=0pt,
  borderline west={2pt}{0pt}{edu-blue!50},
  arc=0pt,
  left=3mm, right=3mm, top=2mm, bottom=2mm,
  before skip=2mm,
  after skip=3mm,
  fontupper=\small,
}

% 解答步骤编号
\newcounter{solstep}
\newcommand{\solstep}[2]{%
  \stepcounter{solstep}%
  \par\medskip
  \noindent\hangindent=6mm
  {\color{edu-blue}\bfseries\textcircled{\scriptsize\thesolstep}\enspace #1}\enspace #2\par
}

\begin{document}

\section*{三垂直与旋转等腰直角专题（二） · 练习卷（教师版）}

\section*{一、选择题（每题 12 分，共 48 分）}


\begin{keyideabox}
  \textbf{一线三垂直变式：直角顶点改变}
  若等腰直角三角形的直角顶点变更为 $B$ （$\angle ABC = 90^\circ$），则必须以点 $B$ 为垂线辅助线的核心端点，从另外两个顶点 $A, C$ 向过 $B$ 的平行线/轴线作垂线。
\end{keyideabox}



\begin{question}[points=12]
  直线 $y = -x + 4$ 与 $x$ 轴、 $y$ 轴交于 $A$ 和 $B$。以 $AB$ 为直角边在第一象限作等腰 Rt$\triangle ABC$ （直角顶点为 $B$，即 $\angle ABC = 90^\circ$），则点 $C$ 的坐标为
  \begin{choices}
    \item $(4, 8)$
    \item $(8, 4)$
    \item $(4, 4)$
    \item $(0, 8)$
  \end{choices}
  \paren[A]
  \begin{solutionblock}
    $A(4,0)$，$B(0,4)$。直角顶点为 $B$，过 $B$ 做水平线（即 $y=4$）。自 $A$ 作垂线垂足 $A_1(4,4)$，$AA_1 = 4$，$BA_1 = 4$。\\ $\triangle ABA_1 \cong \triangle CBB_1$ (AAS) $\implies B_1$ 位于 $B$ 右侧 $4$ 单位，高为 $4$。点 $C$ 坐标为 $(4, 8)$。
  \end{solutionblock}
\end{question}



\begin{question}[points=12]
  点 $A(0, 3)$，点 $B(4, 0)$。以 $AB$ 为底边在第一象限内作等腰 Rt$\triangle ABC$ （直角顶点为 $B$，$\angle ABC = 90^\circ$），则点 $C$ 的坐标为
  \begin{choices}
    \item $(7, 4)$
    \item $(4, 7)$
    \item $(7, 3)$
    \item $(3, 7)$
  \end{choices}
  \paren[A]
  \begin{solutionblock}
    $A(0,3)$，$B(4,0)$。直角顶点为 $B$，过 $B$ 做水平线（$x$轴）。过 $A$ 作垂线，垂足 $A_1(0,0)$。$BA_1 = 4$，$AA_1 = 3$。\\ $\triangle ABA_1 \cong \triangle CBC_1$ (AAS) $\implies CC_1 = BA_1 = 4 \implies y_c = 4$；$BC_1 = AA_1 = 3 \implies x_c = x_b + 3 = 7$。坐标为 $(7, 4)$。
  \end{solutionblock}
\end{question}



\begin{question}[points=12]
  在平面直角坐标系中，等腰 Rt$\triangle ABC$ 的直角顶点在 $B(0, 2)$ 处，锐角顶点 $A(4, 0)$ 在 $x$ 轴上。若点 $C$ 在第一象限内，则点 $C$ 的坐标为
  \begin{choices}
    \item $(2, 6)$
    \item $(6, 2)$
    \item $(2, 4)$
    \item $(4, 6)$
  \end{choices}
  \paren[A]
  \begin{solutionblock}
    $B(0,2)$， $A(4,0)$。过直角顶点 $B$ 沿水平线 $y=2$ 辅助线。过 $A$ 作垂线垂足 $A_1(4,2)$。$BA_1 = 4$，$AA_1 = 2$。\\ $\triangle ABA_1 \cong \triangle CBB_1 \implies B_1$ 在 $B$ 右侧 $2$ 单位，高为 $4$。坐标为 $(2, 6)$。
  \end{solutionblock}
\end{question}



\begin{question}[points=12]
  已知 $A(-2, 0)$，$B(0, 2)$。以 $AB$ 为直角边作等腰 Rt$\triangle ABC$ （直角顶点为 $B$，$\angle ABC = 90^\circ$），点 $C$ 在第一象限，则点 $C$ 的坐标为
  \begin{choices}
    \item $(2, 4)$
    \item $(4, 2)$
    \item $(2, 2)$
    \item $(4, 4)$
  \end{choices}
  \paren[A]
  \begin{solutionblock}
    $A(-2,0)$， $B(0,2)$。直角在 $B$，作水平辅助线。$A_1(-2,2)$， $BA_1 = 2$， $AA_1 = 2$。\\ 全等 $\implies B_1$ 在右侧 $2$，高为 $2$。坐标为 $(2, 4)$。
  \end{solutionblock}
\end{question}


\section*{二、填空题（每题 12 分，共 12 分）}


\begin{question}[points=12]
  已知点 $A(3, 0)$ 和点 $B(0, 1)$。以 $AB$ 为直角边在第一象限作等腰 Rt$\triangle ABC$ （直角顶点为 $B$），则点 $C$ 的坐标为 \fillin。
  \paren[$(1, 4)$]
  \begin{solutionblock}
    $A(3,0)$， $B(0,1)$。过 $B$ 水平线。$A_1(3,1)$。$BA_1 = 3$， $AA_1 = 1$。\\ 全等 $\implies C$ 在 $B$ 右侧 $1$ 个单位，高为 $3$ 个单位（即 $y_c = 1+3=4$）。坐标为 $(1, 4)$。
  \end{solutionblock}
\end{question}


\section*{三、解答题（共 40 分）}

\needspace{8\baselineskip}

\begin{problem}[points=20]
  等腰 Rt$\triangle ABC$ 中，直角顶点位于点 $B(0, 1)$ 处，另一个锐角顶点 $A(\sqrt{3}, 0)$ 在 $x$ 轴上。且点 $C$ 位于第一象限内。

(1) 请画出辅助线，证明一线三垂直对应的两个三角形全等；
(2) 求出点 $C$ 的坐标。
  \answerarea[60mm]
  \setcounter{solstep}{0}
  \begin{solutionblock}
    \textbf{答案：}(1) AAS全等证明；(2) $C(1, \sqrt{3}+1)$\par\medskip
    \solstep{做辅助线与证明全等}{过直角顶点 $B(0,1)$ 做水平辅助线 $y=1$。\\ 过点 $A(\sqrt{3},0)$ 向辅助线作垂线，垂足为 $A_1(\sqrt{3}, 1)$。\\ 过点 $C(x_c, y_c)$ 向辅助线作垂线，垂足为 $C_1(x_c, 1)$。\\ $\because \angle A_1BA + \angle CBC_1 = 90^\circ$ 且 $\angle A_1BA + \angle BAA_1 = 90^\circ \implies \angle BAA_1 = \angle CBC_1$。\\ 结合 $AB = BC$ 与 $\angle AA_1B = \angle BC_1C = 90^\circ$，\\ 判定 $\triangle BAA_1 \cong \triangle CBC_1$ (AAS)。}
    \solstep{通过全等求 C 坐标}{由全等对应边相等：\\ $CC_1 = AA_1 = |0-1| = 1 \implies y_c = 1 + 1 = 2$？ 等等！\\ $AA_1$ 是 $A(\sqrt{3},0)$ 到 $y=1$ 的垂直距离，故 $AA_1 = 1$。\\ $BA_1$ 是水平距离，故 $BA_1 = \sqrt{3}$。\\ 对应关系：$CC_1 = BA_1 = \sqrt{3} \implies y_c = y_b + CC_1 = 1 + \sqrt{3}$。\\ $BC_1 = AA_1 = 1 \implies x_c = x_b + BC_1 = 0 + 1 = 1$。\\ 故点 $C$ 的坐标为 $(1, \sqrt{3}+1)$。}
  \end{solutionblock}
\end{problem}


\needspace{8\baselineskip}

\begin{problem}[points=20]
  已知等腰 Rt$\triangle ABC$ 的直角顶点位于点 $B(0, 2)$，锐角顶点 $A(4, 0)$。

(1) 证明过点 $B$ 作水平线构建的一线三垂直三角形全等；
(2) 求点 $C$ 的坐标，并验证 $BA = BC$。
  \answerarea[55mm]
  \setcounter{solstep}{0}
  \begin{solutionblock}
    \textbf{答案：}(1) AAS全等证明；(2) $C(2, 6)$\par\medskip
    \solstep{全等证明}{过直角顶点 $B(0,2)$ 做水平辅助线 $y=2$。过 $A, C$ 作垂线。\\ 通过互余关系证明 $\triangle BAA_1 \cong \triangle CBC_1$ (AAS)。}
    \solstep{计算 C 坐标}{$AA_1 = 2$，$BA_1 = 4$。\\ 由全等得 $CC_1 = BA_1 = 4 \implies y_c = y_b + 4 = 2+4=6$。\\ $BC_1 = AA_1 = 2 \implies x_c = x_b + 2 = 0+2=2$。\\ 故点 $C$ 坐标为 $(2, 6)$。}
    \solstep{验证边长}{$BA^2 = (4-0)^2 + (0-2)^2 = 16 + 4 = 20 \implies BA = 2\sqrt{5}$。\\ $BC^2 = (2-0)^2 + (6-2)^2 = 4 + 16 = 20 \implies BC = 2\sqrt{5}$。\\ 长度相等，验证通过。}
  \end{solutionblock}
\end{problem}


\clearpage
\section*{参考答案}


\begin{keyideabox}
  \textbf{选择题}
  1. A  2. A  3. A  4. A
\end{keyideabox}



\begin{keyideabox}
  \textbf{填空题}
  1. (1, 4)
\end{keyideabox}



\begin{keyideabox}
  \textbf{解答题}
  第 1 题：(1) AAS全等证明；(2) $C(1, \sqrt{3}+1)$
第 2 题：(1) AAS全等证明；(2) $C(2, 6)$
\end{keyideabox}



\end{document}
